\documentclass[11pt,a4paper]{article}
% ==============================================================
%  Préambule commun aux fiches -- Advanced Probabilities (M1)
%  Usage, dans chaque fiche :
%     \documentclass[11pt,a4paper]{article}
%     % ==============================================================
%  Préambule commun aux fiches -- Advanced Probabilities (M1)
%  Usage, dans chaque fiche :
%     \documentclass[11pt,a4paper]{article}
%     % ==============================================================
%  Préambule commun aux fiches -- Advanced Probabilities (M1)
%  Usage, dans chaque fiche :
%     \documentclass[11pt,a4paper]{article}
%     \input{../preambule.tex}
%     \fiche{C5}{Vecteurs gaussiens}{Gaussian vectors}
% ==============================================================
\usepackage[T1]{fontenc}
\usepackage[utf8]{inputenc}
\usepackage[french]{babel}
\usepackage{lmodern}
\usepackage[a4paper,margin=1.9cm,top=1.7cm,bottom=1.9cm]{geometry}
\usepackage{amsmath,amssymb,amsthm,mathtools}
\usepackage[inline]{enumitem}
\usepackage{booktabs,array,tabularx}
\usepackage{xcolor}
\usepackage[most]{tcolorbox}
\usepackage{fancyhdr}
\usepackage{hyperref}
\hypersetup{colorlinks=true,linkcolor=blue!50!black,urlcolor=blue!50!black}

\setlist{itemsep=1pt,topsep=3pt,parsep=0pt}
\setlength{\parindent}{0pt}
\setlength{\parskip}{3pt}
\emergencystretch=3em

% ---------- Identification de la fiche -------------------------
\newcommand{\ficheid}{}
\newcommand{\fichetitre}{}
\newcommand{\fiche}[3]{%
  \renewcommand{\ficheid}{#1}\renewcommand{\fichetitre}{#2}%
  \pagestyle{fancy}\fancyhf{}%
  \fancyhead[L]{\small\textsf{Advanced Probabilities -- Fiche #1}}%
  \fancyhead[R]{\small\textsf{#2}}%
  \fancyfoot[C]{\small\thepage}%
  \begin{center}
    {\LARGE\bfseries Fiche #1 -- #2}\par\smallskip
    {\large\itshape (#3)}\par\smallskip
    {\small Cours de V.\,M.~Lambert, M1 -- semestre 1}
  \end{center}\vspace{-2pt}\hrule\vspace{6pt}}

% ---------- Environnements numérotés (un seul compteur) --------
\theoremstyle{definition}
\newtheorem{definition}{Définition}
\renewcommand{\thedefinition}{\ficheid.\arabic{definition}}
\newtheorem{theoreme}[definition]{Théorème}
\newtheorem{proposition}[definition]{Proposition}
\newtheorem{lemme}[definition]{Lemme}
\newtheorem{corollaire}[definition]{Corollaire}
\newtheorem{exemple}[definition]{Exemple}
\theoremstyle{remark}
\newtheorem{remarque}[definition]{Remarque}

% Résultat hors cours (standard, signalé comme tel)
\theoremstyle{definition}
\newtheorem{horscours}[definition]{Hors cours}

% Idée de preuve
\newenvironment{idee}{\par\small\textit{Idée de preuve.}~}{\par}

% ---------- Encadrés -----------------------------------------
\tcbset{fichebox/.style={enhanced,breakable,boxrule=0.6pt,arc=2pt,
  left=5pt,right=5pt,top=3pt,bottom=3pt,fonttitle=\bfseries}}
\newtcolorbox{piege}[1][]{fichebox,colback=red!4,colframe=red!60!black,
  title={Pièges classiques},#1}
\newtcolorbox{liens}[1][]{fichebox,colback=blue!3,colframe=blue!50!black,
  title={Liens entre notions},#1}
\newtcolorbox{methode}[1][Méthode]{fichebox,colback=green!4,colframe=green!40!black,title={#1}}
\newtcolorbox{alerte}[1][]{fichebox,colback=orange!6,colframe=orange!70!black,
  title={Attention à l'énoncé},#1}

% Exercice de TD (fiches de rappels)
\newcounter{exo}
\newenvironment{exo}[1][]{\refstepcounter{exo}\par\medskip
  \noindent{\large\bfseries Exercice~\theexo}\ifx\relax#1\relax\else~{\large\bfseries -- #1}\fi\par\nopagebreak\smallskip}{\par}
\newcommand{\outils}{\textbf{Outils du cours.}~}
\newcommand{\meth}{\textbf{Méthode / astuce.}~}
\newcommand{\indic}{\textbf{Indication.}~}
\newcommand{\rep}{\textbf{Réponse.}~}

% ---------- Notations du cours ---------------------------------
\newcommand{\R}{\mathbb{R}}
\newcommand{\N}{\mathbb{N}}
\newcommand{\Q}{\mathbb{Q}}
\newcommand{\C}{\mathbb{C}}
\newcommand{\PP}{\mathbb{P}}
\newcommand{\QQ}{\mathbb{Q}}
\newcommand{\E}{\mathbb{E}}
\newcommand{\F}{\mathcal{F}}
\newcommand{\G}{\mathcal{G}}
\newcommand{\Hc}{\mathcal{H}}
\newcommand{\Tc}{\mathcal{T}}
\newcommand{\Bc}{\mathcal{B}}
\newcommand{\Cc}{\mathcal{C}}
\newcommand{\Ec}{\mathcal{E}}
\newcommand{\Pc}{\mathcal{P}}
\newcommand{\Nc}{\mathcal{N}}
\newcommand{\LN}{\mathcal{LN}}
\newcommand{\Expo}{\mathcal{E}}
\newcommand{\Poi}{\mathcal{P}}
\newcommand{\Ff}{\mathbb{F}}
\newcommand{\1}{\mathbf{1}}
\newcommand{\dd}{\,\mathrm{d}}
\DeclareMathOperator{\Var}{\mathbb{V}}
\DeclareMathOperator{\Cov}{Cov}
\DeclareMathOperator{\card}{card}
\DeclareMathOperator{\diag}{diag}
\newcommand{\indep}{\perp\!\!\!\perp}
\newcommand{\ps}{\ \PP\text{-p.s.}}
\newcommand{\tendloi}{\xrightarrow[n\to+\infty]{\text{loi}}}
\newcommand{\tendps}{\xrightarrow[n\to+\infty]{\text{p.s.}}}
\newcommand{\tendproba}{\xrightarrow[n\to+\infty]{\PP}}
\newcommand{\tendLp}{\xrightarrow[n\to+\infty]{L^p}}
\newcommand{\tend}{\xrightarrow[n\to+\infty]{}}

%     \fiche{C5}{Vecteurs gaussiens}{Gaussian vectors}
% ==============================================================
\usepackage[T1]{fontenc}
\usepackage[utf8]{inputenc}
\usepackage[french]{babel}
\usepackage{lmodern}
\usepackage[a4paper,margin=1.9cm,top=1.7cm,bottom=1.9cm]{geometry}
\usepackage{amsmath,amssymb,amsthm,mathtools}
\usepackage[inline]{enumitem}
\usepackage{booktabs,array,tabularx}
\usepackage{xcolor}
\usepackage[most]{tcolorbox}
\usepackage{fancyhdr}
\usepackage{hyperref}
\hypersetup{colorlinks=true,linkcolor=blue!50!black,urlcolor=blue!50!black}

\setlist{itemsep=1pt,topsep=3pt,parsep=0pt}
\setlength{\parindent}{0pt}
\setlength{\parskip}{3pt}
\emergencystretch=3em

% ---------- Identification de la fiche -------------------------
\newcommand{\ficheid}{}
\newcommand{\fichetitre}{}
\newcommand{\fiche}[3]{%
  \renewcommand{\ficheid}{#1}\renewcommand{\fichetitre}{#2}%
  \pagestyle{fancy}\fancyhf{}%
  \fancyhead[L]{\small\textsf{Advanced Probabilities -- Fiche #1}}%
  \fancyhead[R]{\small\textsf{#2}}%
  \fancyfoot[C]{\small\thepage}%
  \begin{center}
    {\LARGE\bfseries Fiche #1 -- #2}\par\smallskip
    {\large\itshape (#3)}\par\smallskip
    {\small Cours de V.\,M.~Lambert, M1 -- semestre 1}
  \end{center}\vspace{-2pt}\hrule\vspace{6pt}}

% ---------- Environnements numérotés (un seul compteur) --------
\theoremstyle{definition}
\newtheorem{definition}{Définition}
\renewcommand{\thedefinition}{\ficheid.\arabic{definition}}
\newtheorem{theoreme}[definition]{Théorème}
\newtheorem{proposition}[definition]{Proposition}
\newtheorem{lemme}[definition]{Lemme}
\newtheorem{corollaire}[definition]{Corollaire}
\newtheorem{exemple}[definition]{Exemple}
\theoremstyle{remark}
\newtheorem{remarque}[definition]{Remarque}

% Résultat hors cours (standard, signalé comme tel)
\theoremstyle{definition}
\newtheorem{horscours}[definition]{Hors cours}

% Idée de preuve
\newenvironment{idee}{\par\small\textit{Idée de preuve.}~}{\par}

% ---------- Encadrés -----------------------------------------
\tcbset{fichebox/.style={enhanced,breakable,boxrule=0.6pt,arc=2pt,
  left=5pt,right=5pt,top=3pt,bottom=3pt,fonttitle=\bfseries}}
\newtcolorbox{piege}[1][]{fichebox,colback=red!4,colframe=red!60!black,
  title={Pièges classiques},#1}
\newtcolorbox{liens}[1][]{fichebox,colback=blue!3,colframe=blue!50!black,
  title={Liens entre notions},#1}
\newtcolorbox{methode}[1][Méthode]{fichebox,colback=green!4,colframe=green!40!black,title={#1}}
\newtcolorbox{alerte}[1][]{fichebox,colback=orange!6,colframe=orange!70!black,
  title={Attention à l'énoncé},#1}

% Exercice de TD (fiches de rappels)
\newcounter{exo}
\newenvironment{exo}[1][]{\refstepcounter{exo}\par\medskip
  \noindent{\large\bfseries Exercice~\theexo}\ifx\relax#1\relax\else~{\large\bfseries -- #1}\fi\par\nopagebreak\smallskip}{\par}
\newcommand{\outils}{\textbf{Outils du cours.}~}
\newcommand{\meth}{\textbf{Méthode / astuce.}~}
\newcommand{\indic}{\textbf{Indication.}~}
\newcommand{\rep}{\textbf{Réponse.}~}

% ---------- Notations du cours ---------------------------------
\newcommand{\R}{\mathbb{R}}
\newcommand{\N}{\mathbb{N}}
\newcommand{\Q}{\mathbb{Q}}
\newcommand{\C}{\mathbb{C}}
\newcommand{\PP}{\mathbb{P}}
\newcommand{\QQ}{\mathbb{Q}}
\newcommand{\E}{\mathbb{E}}
\newcommand{\F}{\mathcal{F}}
\newcommand{\G}{\mathcal{G}}
\newcommand{\Hc}{\mathcal{H}}
\newcommand{\Tc}{\mathcal{T}}
\newcommand{\Bc}{\mathcal{B}}
\newcommand{\Cc}{\mathcal{C}}
\newcommand{\Ec}{\mathcal{E}}
\newcommand{\Pc}{\mathcal{P}}
\newcommand{\Nc}{\mathcal{N}}
\newcommand{\LN}{\mathcal{LN}}
\newcommand{\Expo}{\mathcal{E}}
\newcommand{\Poi}{\mathcal{P}}
\newcommand{\Ff}{\mathbb{F}}
\newcommand{\1}{\mathbf{1}}
\newcommand{\dd}{\,\mathrm{d}}
\DeclareMathOperator{\Var}{\mathbb{V}}
\DeclareMathOperator{\Cov}{Cov}
\DeclareMathOperator{\card}{card}
\DeclareMathOperator{\diag}{diag}
\newcommand{\indep}{\perp\!\!\!\perp}
\newcommand{\ps}{\ \PP\text{-p.s.}}
\newcommand{\tendloi}{\xrightarrow[n\to+\infty]{\text{loi}}}
\newcommand{\tendps}{\xrightarrow[n\to+\infty]{\text{p.s.}}}
\newcommand{\tendproba}{\xrightarrow[n\to+\infty]{\PP}}
\newcommand{\tendLp}{\xrightarrow[n\to+\infty]{L^p}}
\newcommand{\tend}{\xrightarrow[n\to+\infty]{}}

%     \fiche{C5}{Vecteurs gaussiens}{Gaussian vectors}
% ==============================================================
\usepackage[T1]{fontenc}
\usepackage[utf8]{inputenc}
\usepackage[french]{babel}
\usepackage{lmodern}
\usepackage[a4paper,margin=1.9cm,top=1.7cm,bottom=1.9cm]{geometry}
\usepackage{amsmath,amssymb,amsthm,mathtools}
\usepackage[inline]{enumitem}
\usepackage{booktabs,array,tabularx}
\usepackage{xcolor}
\usepackage[most]{tcolorbox}
\usepackage{fancyhdr}
\usepackage{hyperref}
\hypersetup{colorlinks=true,linkcolor=blue!50!black,urlcolor=blue!50!black}

\setlist{itemsep=1pt,topsep=3pt,parsep=0pt}
\setlength{\parindent}{0pt}
\setlength{\parskip}{3pt}
\emergencystretch=3em

% ---------- Identification de la fiche -------------------------
\newcommand{\ficheid}{}
\newcommand{\fichetitre}{}
\newcommand{\fiche}[3]{%
  \renewcommand{\ficheid}{#1}\renewcommand{\fichetitre}{#2}%
  \pagestyle{fancy}\fancyhf{}%
  \fancyhead[L]{\small\textsf{Advanced Probabilities -- Fiche #1}}%
  \fancyhead[R]{\small\textsf{#2}}%
  \fancyfoot[C]{\small\thepage}%
  \begin{center}
    {\LARGE\bfseries Fiche #1 -- #2}\par\smallskip
    {\large\itshape (#3)}\par\smallskip
    {\small Cours de V.\,M.~Lambert, M1 -- semestre 1}
  \end{center}\vspace{-2pt}\hrule\vspace{6pt}}

% ---------- Environnements numérotés (un seul compteur) --------
\theoremstyle{definition}
\newtheorem{definition}{Définition}
\renewcommand{\thedefinition}{\ficheid.\arabic{definition}}
\newtheorem{theoreme}[definition]{Théorème}
\newtheorem{proposition}[definition]{Proposition}
\newtheorem{lemme}[definition]{Lemme}
\newtheorem{corollaire}[definition]{Corollaire}
\newtheorem{exemple}[definition]{Exemple}
\theoremstyle{remark}
\newtheorem{remarque}[definition]{Remarque}

% Résultat hors cours (standard, signalé comme tel)
\theoremstyle{definition}
\newtheorem{horscours}[definition]{Hors cours}

% Idée de preuve
\newenvironment{idee}{\par\small\textit{Idée de preuve.}~}{\par}

% ---------- Encadrés -----------------------------------------
\tcbset{fichebox/.style={enhanced,breakable,boxrule=0.6pt,arc=2pt,
  left=5pt,right=5pt,top=3pt,bottom=3pt,fonttitle=\bfseries}}
\newtcolorbox{piege}[1][]{fichebox,colback=red!4,colframe=red!60!black,
  title={Pièges classiques},#1}
\newtcolorbox{liens}[1][]{fichebox,colback=blue!3,colframe=blue!50!black,
  title={Liens entre notions},#1}
\newtcolorbox{methode}[1][Méthode]{fichebox,colback=green!4,colframe=green!40!black,title={#1}}
\newtcolorbox{alerte}[1][]{fichebox,colback=orange!6,colframe=orange!70!black,
  title={Attention à l'énoncé},#1}

% Exercice de TD (fiches de rappels)
\newcounter{exo}
\newenvironment{exo}[1][]{\refstepcounter{exo}\par\medskip
  \noindent{\large\bfseries Exercice~\theexo}\ifx\relax#1\relax\else~{\large\bfseries -- #1}\fi\par\nopagebreak\smallskip}{\par}
\newcommand{\outils}{\textbf{Outils du cours.}~}
\newcommand{\meth}{\textbf{Méthode / astuce.}~}
\newcommand{\indic}{\textbf{Indication.}~}
\newcommand{\rep}{\textbf{Réponse.}~}

% ---------- Notations du cours ---------------------------------
\newcommand{\R}{\mathbb{R}}
\newcommand{\N}{\mathbb{N}}
\newcommand{\Q}{\mathbb{Q}}
\newcommand{\C}{\mathbb{C}}
\newcommand{\PP}{\mathbb{P}}
\newcommand{\QQ}{\mathbb{Q}}
\newcommand{\E}{\mathbb{E}}
\newcommand{\F}{\mathcal{F}}
\newcommand{\G}{\mathcal{G}}
\newcommand{\Hc}{\mathcal{H}}
\newcommand{\Tc}{\mathcal{T}}
\newcommand{\Bc}{\mathcal{B}}
\newcommand{\Cc}{\mathcal{C}}
\newcommand{\Ec}{\mathcal{E}}
\newcommand{\Pc}{\mathcal{P}}
\newcommand{\Nc}{\mathcal{N}}
\newcommand{\LN}{\mathcal{LN}}
\newcommand{\Expo}{\mathcal{E}}
\newcommand{\Poi}{\mathcal{P}}
\newcommand{\Ff}{\mathbb{F}}
\newcommand{\1}{\mathbf{1}}
\newcommand{\dd}{\,\mathrm{d}}
\DeclareMathOperator{\Var}{\mathbb{V}}
\DeclareMathOperator{\Cov}{Cov}
\DeclareMathOperator{\card}{card}
\DeclareMathOperator{\diag}{diag}
\newcommand{\indep}{\perp\!\!\!\perp}
\newcommand{\ps}{\ \PP\text{-p.s.}}
\newcommand{\tendloi}{\xrightarrow[n\to+\infty]{\text{loi}}}
\newcommand{\tendps}{\xrightarrow[n\to+\infty]{\text{p.s.}}}
\newcommand{\tendproba}{\xrightarrow[n\to+\infty]{\PP}}
\newcommand{\tendLp}{\xrightarrow[n\to+\infty]{L^p}}
\newcommand{\tend}{\xrightarrow[n\to+\infty]{}}

\begin{document}
\fiche{TD6}{Martingales et théorème d'arrêt}{Doob's stopping time theorem}

\textit{À lire avant de chercher le TD. Renvois : « C8.11 » = fiche C8, énoncé n°11. Fiche associée : \textbf{C8} (et C6 pour les
calculs conditionnels). Les outils « martingale régulière » et « passage à la limite dans l'arrêt » sont \textbf{hors cours} (C8.11).}

\begin{exo}[{Martingale exponentielle et sortie de $]-a,a[$}]
\begin{alerte}
L'énoncé ne dit pas explicitement que les $Z_i$ sont \textbf{indépendantes} (ni qu'elles sont définies pour tout $i\ge1$) : il faut le supposer,
c'est le cadre de la marche aléatoire simple. « $\F_0=(\Omega,\emptyset)$ » désigne la tribu triviale $\{\emptyset,\Omega\}$.
% Vérifié sur le PDF de l'énoncé (p. 1) : aucune hypothèse d'indépendance des Z_i.
\end{alerte}
\outils Martingale en temps discret (C8.4) ; sortir le connu, indépendance (C6.5) ; martingale arrêtée et $\E(X_{n\wedge\tau})=\E(X_0)$ (C8.10, C8.11) ;
convergence monotone et dominée ; martingale régulière (C8.11).

\meth
\begin{enumerate}
  \item $\cos(\lambda S_{n+1})=\cos(\lambda S_n)\cos(\lambda Z_{n+1})-\sin(\lambda S_n)\sin(\lambda Z_{n+1})$ ; or $\E\cos(\lambda Z)=\cos\lambda$ et
  $\E\sin(\lambda Z)=0$. Noter que $\cos\lambda>0$ car $0<\lambda<\frac\pi{2a}\le\frac\pi2$.
  \item $n\wedge\tau$ est un temps d'arrêt borné : $\E(X_{n\wedge\tau})=X_0=1$. Avant la sortie, $|S_{n\wedge\tau}|\le a$, donc
  $\lambda|S_{n\wedge\tau}|\le\lambda a<\frac\pi2$ et $\cos(\lambda S_{n\wedge\tau})\ge\cos(\lambda a)>0$.
  \item Convergence monotone : $(\cos\lambda)^{-n\wedge\tau}\uparrow(\cos\lambda)^{-\tau}$, qui vaut $+\infty$ sur $\{\tau=\infty\}$ (car $0<\cos\lambda<1$).
  \item Domination : $|X_{n\wedge\tau}|\le(\cos\lambda)^{-\tau}\in L^1$, donc la martingale arrêtée est uniformément intégrable (régulière).
  \item Passer à la limite dans $\E(X_{n\wedge\tau})=1$ (convergence dominée) et utiliser $|S_\tau|=a$, donc $\cos(\lambda S_\tau)=\cos(\lambda a)$.
  \item $(\cos\lambda)^{-\tau}=e^{c\tau}$ avec $c=-\ln\cos\lambda>0$ : un moment exponentiel fini domine tous les moments polynomiaux
  ($\tau^p\le C_pe^{c\tau}$).
\end{enumerate}

\rep 5. $\E\bigl((\cos\lambda)^{-\tau}\bigr)=\dfrac1{\cos(\lambda a)}$. 6. Oui, $\tau\in L^p$ pour tout $p\ge1$.
\end{exo}

\begin{exo}[Ruine de l'investisseur]
\begin{alerte}
\begin{itemize}
  \item $X_{n+1}=X_n\pm10$ avec probabilités $p$ et $1-p$ donne $\E(X_{n+1}\mid\F_n)=X_n+10(2p-1)$ sur $\{n<T\}$. Le processus (arrêté)
  n'est une \textbf{martingale que si $p=\frac12$} ; c'est une sous-martingale si $p>\frac12$, une sur-martingale si $p<\frac12$.
  \item L'énoncé affirme que $X_T$ vaut $150$ avec probabilité $p$ : c'est \textbf{faux en général}. Pour $p=\frac12$, l'arrêt de Doob donne
  $\PP(X_T=150)=\frac{100}{150}=\frac23$ (cf.\ exemple C8.13 avec $a=10$ et $b=5$ pas de $10$).
  \item Les questions 6--10 reprennent les questions 1--5.
\end{itemize}
% Vérifié : E(X_{n+1}|F_n) = X_n + 10(2p-1), martingale ssi p = 1/2, et alors P(X_T=150) = 2/3.
\end{alerte}
\outils Martingale (C8.2, C8.4) ; arrêt de Doob (C8.10) ; passage à la limite pour $T$ non borné (C8.11) ; exemple C8.13 ; probabilités totales (C1.13).

\meth
\begin{itemize}
  \item Q1/Q6 : probabilités totales sur le résultat de la période $n+1$.
  \item Q2/Q7 : $(X_{n\wedge T})$ est bornée ($0\le X_{n\wedge T}\le150$) ; montrer $T<\infty$ p.s.\ (par ex.\ avec la martingale $M_n=(X_n-100)^2-100n$
  comme dans C8.13 si $p=\frac12$), puis convergence dominée dans $\E(X_{n\wedge T})=100$.
  \item Cas $p\neq\frac12$ (hors cours, ruine du joueur) : $\bigl(\frac{1-p}p\bigr)^{X_{n}/10}$ est une martingale ; même schéma d'arrêt.
  \item Q3/Q8 : $\E(X_T)=150\,\PP(X_T=150)+0$.
\end{itemize}

\rep Pour $p=\frac12$ : $\PP(X_T=150)=\frac23$ et $\E(X_T)=100=X_0$ (le jeu équitable ne crée pas de gain moyen).
Pour $p\neq\frac12$, avec $r=\frac{1-p}{p}$ (départ à 10 pas de 0, cible à 15 pas de 0, pas $=10$) : $\PP(X_T=A)=\dfrac{1-r^{10}}{1-r^{15}}$ et $\E(X_T)=A\,\PP(X_T=A)$, $A=150$.
Exemples : $p=0{,}6$ donne $\PP\approx0{,}985$ et $\E(X_T)\approx147{,}7$ ; $p=0{,}4$ donne $\PP\approx0{,}130$ et $\E(X_T)\approx19{,}5$.
Interprétation (Q4, 5, 9, 10) : la probabilité d'atteindre l'objectif croît avec $p$ ; si $p<\frac12$ l'espérance de gain est négative et
la ruine est probable ; un seuil $A$ plus élevé diminue la probabilité de l'atteindre.
\end{exo}

\begin{exo}[Sortie d'un intervalle asymétrique]
\outils Comme à l'exercice 1 (C8.4, C8.10, C8.11).

\meth 1. Même calcul qu'à l'exercice 1 Q1 avec $x=\lambda(S_n-\frac{b-a}2)$ : $\E\cos(x+\lambda Z)=\cos x\cos\lambda$.
2. Sur $[-a,b]$, $|S_{n\wedge T}-\frac{b-a}2|\le\frac{a+b}2$ ; pour $0<\lambda<\frac\pi{a+b}$, $\cos\bigl(\lambda(S_{n\wedge T}-\frac{b-a}2)\bigr)\ge\cos\frac{\lambda(a+b)}2>0$.
Reprendre les étapes 2 à 5 de l'exercice 1. En $S_T=b$ comme en $S_T=-a$, $\bigl|S_T-\frac{b-a}2\bigr|=\frac{a+b}2$ (cosinus pair).
\end{exo}

\section*{Synthèse : type de question $\to$ méthode}
\renewcommand{\arraystretch}{1.3}
\begin{tabularx}{\linewidth}{@{}>{\raggedright\arraybackslash}p{0.36\linewidth}X@{}}
\toprule
Type de question & Méthode\\\midrule
Martingale trigonométrique / exponentielle & Formule d'addition + indépendance du pas : $\E\cos(x+\lambda Z)=\cos x\cos\lambda$.\\
Appliquer l'arrêt avec $\tau$ non borné & Arrêter en $n\wedge\tau$ (borné, C8.10), puis $n\to\infty$ (C8.11).\\
Montrer $\tau<\infty$ p.s. & Borne uniforme sur $\E(f(n\wedge\tau))$ avec $f\uparrow+\infty$, puis convergence monotone.\\
Passer à la limite dans $\E(X_{n\wedge\tau})$ & Convergence dominée (processus arrêté borné ou dominé par une v.a.\ intégrable).\\
Martingale régulière & Domination par une v.a.\ $L^1$ $\Rightarrow$ uniformément intégrable (C8.11).\\
Moments de $\tau$ & Un moment exponentiel fini entraîne tous les moments polynomiaux.\\
Probabilité de sortie par un bord & $\E(X_\tau)=\E(X_0)$ avec $X_\tau$ à deux valeurs : équation linéaire en $\PP(\text{bord})$.\\
Jeu biaisé ($p\neq\frac12$) & Chercher une martingale $r^{S_n}$ avec $r=\frac{1-p}p$ (hors cours).\\\bottomrule
\end{tabularx}

\end{document}
